\section{Introduction}

\citet{lusardiMitchell2014} and a large literature document that financially literate households make systematically better portfolio decisions: they diversify more, panic less, time the market less, and pay more attention to fees and volatility-adjusted risk. The variation across U.S. states is substantial --- roughly 28 percent of adults answer three of five basic literacy questions correctly in high-literacy states like Minnesota and New Hampshire, against around 21 percent in Mississippi and West Virginia. The local-bias literature supplies the second piece: retail investors hold portfolios tilted toward locally headquartered firms \citep{covalMoskowitz1999, ivkovicWeisbenner2005}, and this local demand moves prices \citep{hongKubikStein2008, bernileKumarSulaeman2015}. Putting the two together yields a clean hypothesis: where retail investors are less literate, they should trade more like the textbook noise trader, so state literacy should moderate the momentum--idiosyncratic-volatility (IV) interaction --- a retail-friendly anomaly combined with the standard limits-to-arbitrage measure --- with the effect weaker in literate states.

In aggregate, however, the data do not deliver. Regressing the three-way interaction on the full panel of CRSP-listed common stock returns a modest negative coefficient with marginal significance, and the magnitude collapses by roughly forty fold once headquarters-state assignment is corrected for the 14 percent of firm-months whose static Compustat snapshot does not match their then-current 10-K filing. The aggregate test is misleading in two ways: it pools firms with very different ownership structures, and it pools firms whose local-investor base is well-defined (stable headquarters) with firms whose base is in transition (relocators). The aggregate null is therefore not informative about whether a literacy channel exists --- only about whether it operates uniformly across the cross-section. It does not.

The contribution of this paper is to ask a sharper question: where in the cross-section does the literacy mechanism become visible, under what conditions on firm ownership and size, and what does the answer identify about the microfoundation? On the stable-headquarters subsample of CRSP-listed common stock from 2009 to 2023, the literacy moderation concentrates in firms whose institutional ownership sits in the middle of the within-stock IO distribution --- the mid-IO tercile, with mean ownership share around fifty percent. Firms with predominantly retail ownership (low IO) show essentially no moderation; firms with predominantly institutional ownership (high IO) also show none. The mid-tercile coefficient survives the addition of firm-size moderation as a control, so the finding is not a size effect dressed in ownership clothing even though mid-IO firms are predominantly mid-cap. Within mid-IO, the coefficient corresponds to a cross-state literacy gradient in exactly the predicted direction: higher-literacy states show weaker anomaly slopes, and the across-state correlation between per-state coefficients and state literacy is sharp and clean.

The economic interpretation the data identify is an \emph{interaction zone}. The literacy channel requires two ingredients to leave a price footprint: retail demand large enough to move prices, and institutional activity sufficient to translate that demand into return differences detectable across states. Where neither is present (micro-cap firms held mostly by retail with no institutional counterparty), retail mispricing is genuine but accumulates uniformly without loading on cross-state literacy variation. Where institutions dominate (large firms with high IO), arbitrage absorbs retail demand and neutralizes the local-state channel. The interaction zone --- firms with mixed ownership at meaningful market size --- is where both ingredients are present. In the full stable-HQ panel it falls predominantly in mid-IO; dropping the bottom 50 percent of firms by market capitalization shifts the moderation to low-IO, and the cross-state gradient migrates with it. It is the same economic population, relocated to a different tercile of the conditional IO distribution.

This migration is the central piece of mechanism evidence the paper introduces --- the empirical signature an interaction-zone reading predicts and the IO-tercile result alone does not deliver. The high-IO tercile shows no cross-state gradient under any sample restriction, a clean predicted-null check the reading passes. The IV-decile decomposition within mid-IO, which would test the canonical limits-to-arbitrage microfoundation under which the effect strengthens at high IV, fails to reject decile equality and returns a near-zero coefficient at the highest IV decile. The textbook limits-to-arbitrage reading is rejected; the interaction-zone reading, identified through size and ownership conditioning rather than through IV, is identified positively by the migration.

Three bounds discipline the result. First, the continuous-margin quadratic-IO U-shape clears five percent on a quadratic specification only --- not under log, kink, or twenty-bin alternatives --- so the functional-form-free claim is the discrete-tercile contrast. Second, under point-in-time literacy reassignment on the full panel the mid-IO coefficient collapses by roughly twenty-six fold; the characterization is a statement about firms with stable headquarters, not the corrected full panel. Third, an instrument using state K-12 personal-finance graduation requirements fails its first-stage relevance test in the right direction (mandate states are systematically less literate, having adopted mandates as a remedy), so the OLS estimate is the preferred inferential statement and the IV failure is documented honestly.

The closest prior work is \citet{korniotisKumar2013}, who show that state-level business cycles predict returns of locally headquartered firms. This paper differs on three margins: the state variable is NFCS literacy rather than the business cycle, the object is a three-way moderation of an anomaly interaction rather than a level prediction, and the discriminating margin is a within-stock IO contrast the business-cycle channel does not predict. A business-cycle control block attenuates the literacy three-way by 20--23 percent, so the two state moderators are distinct but partly overlapping. \citet{korniotisKumarPage2020} use brokerage-demographic sophistication proxies but neither NFCS literacy nor the momentum-IV interaction; \citet{hillertJacobsMuller2014} is the closest precedent for a state-level moderator of momentum, using media coverage with state individualism as the cross-sectional moderator rather than literacy.

\paragraph{Roadmap.} Section~\ref{sec:data} describes the data, the NFCS state-literacy panel, and the stable-headquarters subsample. Section~\ref{sec:identification} states the specification and identifying assumptions. Section~\ref{sec:results} reports the mid-IO concentration, the size-controlled extension, the migration, and the cross-state gradient. Section~\ref{sec:mechanism} develops the interaction-zone characterization. Section~\ref{sec:robustness} reports the disciplining battery. Section~\ref{sec:discussion} discusses interpretation and limitations, and Section~\ref{sec:conclusion} concludes.
